\section*{Nudel-Schinken-Auflauf}
\ihead{}\chead{Personen:4}\ohead{}
\ifoot{Vorbereitungszeit:20 min}\cfoot{}\ofoot{Kochzeit:40 min}
{\Large Zutaten}
\begin{multicols}{2}
\begin{itemize}
\item \SI{500}{g} Fusilli, roh
\item \SI{500}{g} Kochschinken
\item \SI{500}{ml} Sahne
\item \SI{700}{ml} Milch
\item \SI{50}{g} Butter
\item \num{2} Eigelb
\item \num{2} große Zwiebeln
\item \SI{2}{EL} Mehl
\item \num{2} Handvoll frische Petersilie
\item Parmesan, gerieben, nach Belieben
\item Muskat nach Belieben
\item schwarzer Pfeffer nach Belieben
\item \SI{400}{g} Käse, gerieben (z.,B. Gouda oder Emmentaler)
\item \SI{2}{EL} Butter für die Form
\end{itemize}
% \columnbreak
\end{multicols}
\noindent {\Large Zubereitung}\ 
\textit{Vorbereitung:} Backofen auf \SI{200}{\celsius} Ober-/Unterhitze vorheizen.
Zwiebeln und Schinken würfeln und in der Butter anbraten.
Das Mehl hinzugeben und unter Rühren leicht bräunen lassen.
Mit Sahne und Milch ablöschen und aufkochen lassen.
Vom Herd nehmen und die Eigelbe unter Rühren hinzugeben.
Mit Parmesan, Petersilie, Pfeffer und Muskat abschmecken.
Die ungekochten Nudeln in eine gefettete Auflaufform geben und mit der Sauce vermischen.
Den geriebenen Käse darüberstreuen.
Bei \SI{200}{\celsius} Ober-/Unterhitze ca. \SI{30}{min} im vorgeheizten Backofen backen, bis der Käse goldbraun ist.
